\documentclass[12pt,a4paper]{amsart}

\usepackage{amsmath,amssymb,amsthm,mathtools,enumitem,booktabs,array,hyperref}
\usepackage[utf8]{inputenc}
\usepackage[T1]{fontenc}
\usepackage{geometry}
\geometry{margin=2.5cm}

% -------------------------------------------------------
% Theorem environments
% -------------------------------------------------------
\newtheorem{theorem}{Theorem}[section]
\newtheorem{lemma}[theorem]{Lemma}
\newtheorem{corollary}[theorem]{Corollary}
\newtheorem{proposition}[theorem]{Proposition}
\theoremstyle{definition}
\newtheorem{definition}[theorem]{Definition}
\newtheorem{example}[theorem]{Example}
\newtheorem{remark}[theorem]{Remark}
\newtheorem{conjecture}[theorem]{Conjecture}

% -------------------------------------------------------
% Custom commands
% -------------------------------------------------------
\newcommand{\N}{\mathbb{N}}
\newcommand{\Z}{\mathbb{Z}}
\newcommand{\Q}{\mathbb{Q}}
\newcommand{\R}{\mathbb{R}}
\newcommand{\C}{\mathbb{C}}
\newcommand{\F}{\mathbb{F}}
\newcommand{\Fp}{\mathbb{F}_p}
\newcommand{\A}{\mathbb{A}}
\newcommand{\PP}{\mathbb{P}}
\newcommand{\Gal}{\mathrm{Gal}}
\newcommand{\End}{\mathrm{End}}
\newcommand{\Aut}{\mathrm{Aut}}
\newcommand{\id}{\mathrm{id}}
\newcommand{\ch}{\mathrm{ch}}
\newcommand{\td}{\mathrm{td}}
\newcommand{\NS}{\mathrm{NS}}
\newcommand{\Pic}{\mathrm{Pic}}
\newcommand{\CH}{\mathrm{CH}}
\newcommand{\cl}{\mathrm{cl}}
\newcommand{\Num}{\mathrm{Num}}
\DeclareMathOperator{\rank}{rank}
\DeclareMathOperator{\tr}{tr}
\DeclareMathOperator{\codim}{codim}
\DeclareMathOperator{\Spec}{Spec}
\DeclareMathOperator{\Proj}{Proj}

% -------------------------------------------------------
\begin{document}
% -------------------------------------------------------

\title{Grothendieck's Standard Conjectures on Algebraic Cycles:\\
       Weil Cohomology, Motives, and the Lefschetz Operator}

\author{Michael Fuhrmann}
\address{Specialist Mathematics Project, Build 170}
\email{specialist-maths@localhost}
\date{\today}

\begin{abstract}
Grothendieck's four standard conjectures --- Conjecture~A (Lefschetz),
Conjecture~B (Hodge standard), Conjecture~C (K\"unneth), and Conjecture~D
(numerical equals homological) --- form the foundational pillars of the
theory of motives.  Formulated in the late 1960s as algebraic analogues of
classical hard Lefschetz and Hodge-theoretic facts over $\mathbb{C}$, these
conjectures remain almost entirely open for varieties over fields of positive
characteristic, and even over number fields.  In this paper we give precise
and complete statements of all four conjectures within the framework of Weil
cohomology theories, explain their mutual implications, relate them to the
Weil conjectures (proved by Deligne in 1974), the Hodge conjecture, and the
Tate conjecture, and survey the partial results known: Kleiman's proof that
$\mathbf{D}$ implies $\mathbf{A}$, Jannsen's semisimplicity theorem for
numerical motives, Andr\'e's theory of motivated cycles, and the current
state of the art.  We also discuss why these conjectures are so hard and
what a proof would imply for the category of pure motives.
\end{abstract}

\maketitle

\tableofcontents

% -------------------------------------------------------
\section{Introduction}
% -------------------------------------------------------

The central aim of algebraic geometry is to understand the geometry of
algebraic varieties through cohomological invariants.  Over the complex
numbers this programme reaches its apex in classical Hodge theory: the
hard Lefschetz theorem, the Hodge decomposition, and the positive-definiteness
of the Weil form give a rich structure to $H^*(X,\Q)$ for any smooth
projective complex manifold~$X$.

Grothendieck's visionary insight, announced at the 1968 International Congress
of Mathematicians and elaborated in his letter to Serre and other manuscripts,
was that the deepest properties of algebraic cycles --- not just over
$\mathbb{C}$ but over \emph{any} field $k$ --- should be captured by a small
number of precise conjectures that parallel classical Hodge theory.
The four \emph{standard conjectures} bear the labels A, B, C, D (sometimes
also I, II for A and B in Grothendieck's original notation).

These conjectures serve a double purpose.  First, taken together they imply
that the na\"{\i}ve category of correspondences modulo homological equivalence
is actually a semi-simple abelian category --- the category of \emph{pure
motives} $\mathcal{M}_{\mathrm{hom}}(k)$ --- and that this category is
independent of the choice of Weil cohomology theory.  Second, they give
cohomological content to the Weil conjectures (which Deligne proved without
using them by a completely different method).

\subsection*{Organisation of the paper}

Section~\ref{sec:weil} recalls the axioms of a Weil cohomology theory.
Section~\ref{sec:cycles} reviews algebraic cycles and their equivalence
relations.  Section~\ref{sec:lefschetz} introduces the Lefschetz operator
and the classical hard Lefschetz theorem over $\mathbb{C}$.
Sections~\ref{sec:conjA}--\ref{sec:conjD} give complete, rigorous statements
of Conjectures~A--D.  Section~\ref{sec:implications} works out the
logical dependencies among the four conjectures.
Section~\ref{sec:weilconj} relates the standard conjectures to the Weil
conjectures and Deligne's theorem.  Section~\ref{sec:hodge} discusses the
relationship with the Hodge and Tate conjectures.
Section~\ref{sec:motives} explains the motivic consequences.
Section~\ref{sec:andre} surveys Andr\'e motivated cycles.
Section~\ref{sec:jannsen} presents Jannsen's semisimplicity theorem.
Section~\ref{sec:known} collects all known special cases.
Section~\ref{sec:obstacles} explains the main obstacles to a general proof.
Section~\ref{sec:outlook} gives an outlook.

% -------------------------------------------------------
\section{Weil Cohomology Theories}\label{sec:weil}
% -------------------------------------------------------

Throughout, $k$ is a field (often a finite field $\Fp$, a number field, or
$\C$), and all varieties are smooth projective and geometrically connected
over $k$ unless otherwise stated.

\begin{definition}[Weil cohomology theory]
A \emph{Weil cohomology theory with coefficient field $F$ (of characteristic
zero)} is a contravariant functor $H^*$ from the category of smooth projective
$k$-varieties to graded $F$-algebras, together with a cycle class map,
satisfying the following axioms:
\begin{enumerate}[label=(\roman*)]
  \item \textbf{Finite-dimensionality.}  $H^i(X) = 0$ for $i < 0$ or
        $i > 2\dim X$, and each $H^i(X)$ is a finite-dimensional $F$-vector
        space.
  \item \textbf{Poincar\'e duality.}  There is a canonical isomorphism
        $H^{2d}(X)(d) \cong F$ (the \emph{fundamental class}), and the
        pairing $H^i(X) \otimes H^{2d-i}(X) \to F$ induced by cup product
        is a perfect duality.  Here $d = \dim X$.
  \item \textbf{K\"unneth formula.}  There is a natural isomorphism
        $H^*(X \times Y) \cong H^*(X) \otimes_F H^*(Y)$.
  \item \textbf{Cycle class map.}  For each $i$, there is a
        $\Q$-linear homomorphism
        $\cl: \CH^i(X) \to H^{2i}(X)(i)$
        from the Chow group of codimension-$i$ cycles, compatible with
        proper push-forward, flat pull-back, and cup product.
  \item \textbf{Weak Lefschetz / hyperplane section.}  If $Y \hookrightarrow X$
        is a smooth hyperplane section, then
        $H^i(X) \to H^i(Y)$ is an isomorphism for $i < \dim Y$ and
        injective for $i = \dim Y$.
\end{enumerate}
\end{definition}

\begin{example}[Standard examples of Weil cohomology theories]
\begin{enumerate}[label=(\alph*)]
  \item \textbf{Singular cohomology.}  For $k \hookrightarrow \C$, take
        $H^i(X) = H^i(X(\C),\Q)$ with $F = \Q$.
  \item \textbf{de Rham cohomology.}  For $k$ of characteristic zero,
        $H^i_{\mathrm{dR}}(X/k)$ with $F = k$.
  \item \textbf{$\ell$-adic cohomology.}  For $\ell \neq \ch(k)$ (prime),
        $H^i_{\mathrm{\acute{e}t}}(X_{\bar k}, \Q_\ell)$ with $F = \Q_\ell$.
  \item \textbf{Crystalline cohomology.}  For $k = \Fp$,
        $H^i_{\mathrm{cris}}(X/W(k))[\frac{1}{p}]$ with $F = W(k)[\frac{1}{p}]$,
        due to Berthelot.
\end{enumerate}
\end{example}

\begin{remark}
All known Weil cohomology theories satisfy the axioms above (the K\"unneth
formula for $\ell$-adic cohomology is a non-trivial theorem).  A key feature
is that they all give the \emph{same} Betti numbers $b_i(X) = \dim_F H^i(X)$
when they are defined simultaneously, but proving that the full cohomology
algebras are isomorphic requires the standard conjectures.
\end{remark}

% -------------------------------------------------------
\section{Algebraic Cycles and Equivalence Relations}\label{sec:cycles}
% -------------------------------------------------------

Let $X$ be a smooth projective variety of dimension $d$ over $k$.

\begin{definition}[Algebraic cycle]
An \emph{algebraic cycle} of codimension $p$ on $X$ is a finite formal
$\Z$-linear combination $Z = \sum n_i [V_i]$ of irreducible closed
subvarieties $V_i \subset X$ of codimension $p$.  The group of such cycles
is denoted $\mathcal{Z}^p(X)$.
\end{definition}

There are several important equivalence relations on $\mathcal{Z}^p(X)$:

\begin{definition}[Equivalence relations on cycles]
\begin{enumerate}[label=(\roman*)]
  \item \textbf{Rational equivalence}: $Z \sim_{\mathrm{rat}} 0$ if $Z$ is
        the divisor of a rational function on some subvariety.  The quotient
        $\CH^p(X) = \mathcal{Z}^p(X)/{\sim_{\mathrm{rat}}}$ is the
        \emph{Chow group}.
  \item \textbf{Algebraic equivalence}: $Z \sim_{\mathrm{alg}} 0$ if $Z$ is
        algebraically equivalent to zero (parametrised by a connected curve).
  \item \textbf{Homological equivalence}: $Z \sim_{\mathrm{hom}} 0$ if
        $\cl(Z) = 0$ in $H^{2p}(X)(p)$ for some (equivalently, by Conjecture~D,
        all) Weil cohomology theories.
  \item \textbf{Numerical equivalence}: $Z \sim_{\mathrm{num}} 0$ if
        $(Z \cdot W) = 0$ for all cycles $W$ of complementary codimension.
\end{enumerate}
We have the inclusions
${\sim_{\mathrm{rat}}} \supseteq {\sim_{\mathrm{alg}}} \supseteq
 {\sim_{\mathrm{hom}}} \supseteq {\sim_{\mathrm{num}}}$.
\end{definition}

\begin{remark}
The first inclusion $\supseteq$ reflects that rationally equivalent cycles
are a fortiori algebraically equivalent (take a $\PP^1$-parameter).
The last inclusion follows because the intersection form factors through the
cycle class map.  That ${\sim_{\mathrm{hom}}} = {\sim_{\mathrm{num}}}$
is precisely Conjecture~D.
\end{remark}

\begin{definition}[Correspondences]
A \emph{correspondence} from $X$ to $Y$ (of degree $r$) is an element of
$\CH^{\dim X + r}(X \times Y)$.  The category of \emph{pure motives} modulo
an adequate equivalence relation $\sim$ has smooth projective varieties as
objects, with morphisms given by correspondences modulo $\sim$, composed via
$[Z] \circ [W] = (p_{13})_* (p_{12}^* [Z] \cdot p_{23}^* [W])$.
\end{definition}

% -------------------------------------------------------
\section{The Lefschetz Operator and Hard Lefschetz}\label{sec:lefschetz}
% -------------------------------------------------------

Fix a smooth projective variety $X$ of dimension $d$ over $k$, and let
$\xi \in \CH^1(X)$ be the class of a hyperplane section (i.e.\ a very ample
line bundle $L$).  Its cycle class is $L = \cl(\xi) \in H^2(X)(1)$.

\begin{definition}[Lefschetz operator]
The \emph{Lefschetz operator} is the map
\[
  \mathbf{L} = \cup L : H^i(X) \longrightarrow H^{i+2}(X)(1),
  \quad
  \alpha \mapsto L \cup \alpha.
\]
Its iterated powers are $\mathbf{L}^r : H^i(X) \to H^{i+2r}(X)(r)$.
\end{definition}

\begin{theorem}[Hard Lefschetz over $\mathbb{C}$, Hodge 1950]\label{thm:HL}
Let $X$ be a smooth projective complex variety of dimension $d$, and let
$L \in H^2(X,\Q)$ be the class of a hyperplane section.  Then for all
$0 \leq i \leq d$:
\[
  \mathbf{L}^{d-i} : H^i(X,\Q) \xrightarrow{\;\sim\;} H^{2d-i}(X,\Q)
\]
is an isomorphism.
\end{theorem}

This theorem, while classical over $\C$, is \emph{unknown} over general
fields.  Its algebraic analogue is the content of Conjecture~B.

\begin{definition}[Primitive cohomology and Lefschetz decomposition]
For $i \leq d$, the \emph{primitive cohomology} is
$H^i_{\mathrm{prim}}(X) = \ker(\mathbf{L}^{d-i+1}: H^i(X) \to H^{2d-i+2}(X))$.
The Lefschetz decomposition (over $\C$) states:
\[
  H^i(X) = \bigoplus_{r \geq 0} \mathbf{L}^r H^{i-2r}_{\mathrm{prim}}(X).
\]
\end{definition}

\begin{definition}[Lefschetz involution $\Lambda$]
Define the \emph{Lefschetz involution} (or \emph{Lefschetz star operator})
$\Lambda: H^i(X) \to H^{i-2}(X)$ as the adjoint of $\mathbf{L}$ with
respect to the Poincar\'e pairing.  More precisely, for $i \geq d$,
hard Lefschetz provides an isomorphism $\mathbf{L}^{i-d}: H^{2d-i}(X) \to
H^i(X)$, and $\Lambda = (\mathbf{L}^{i-d})^{-1} \circ$ (something).
See Section~\ref{sec:conjB} for the precise algebraic formulation.
\end{definition}

% -------------------------------------------------------
\section{Conjecture~A: The Lefschetz Conjecture}\label{sec:conjA}
% -------------------------------------------------------

Conjecture~A asserts that the operator $\mathbf{L}^{d-i}$ is ``algebraic''
in the sense that it is induced by an algebraic cycle (a correspondence).

\begin{conjecture}[Standard Conjecture~A, Grothendieck 1969]\label{conj:A}
Let $X$ be a smooth projective variety of dimension $d$ over a field $k$,
and let $H^*$ be a Weil cohomology theory with coefficient field $F$.
For each $0 \leq i \leq d$, the isomorphism
\[
  \mathbf{L}^{d-i} : H^i(X) \xrightarrow{\;\sim\;} H^{2d-i}(X)(d-i)
\]
given by iterated cup product with $L = \cl(\xi)$ is algebraic, i.e.\ it is
induced by an algebraic correspondence $\Pi_i \in \CH^d(X \times X)$ (up to
the twist).  Equivalently, the projectors $\pi_i$ in the K\"unneth decomposition
of $[\Delta_X] = \sum_{i=0}^{2d} \pi_i$ are algebraic for all $i$.
\end{conjecture}

\begin{remark}
Conjecture~A has two equivalent formulations:
\begin{enumerate}[label=(\alph*)]
  \item The hard Lefschetz isomorphism $\mathbf{L}^{d-i}$ is algebraic.
  \item The K\"unneth components $\pi_i$ of the diagonal class are algebraic.
\end{enumerate}
The equivalence $(a) \Leftrightarrow (b)$ is a formal consequence of the
decomposition of the identity; see Kleiman~\cite{kleiman1968}.
\end{remark}

\begin{theorem}[Grothendieck]\label{thm:A_implies_C}
Conjecture~A implies Conjecture~C.
\end{theorem}

\begin{proof}[Proof sketch]
Conjecture~A says that $\mathbf{L}^{d-i}$ is algebraic, which gives an
algebraic cycle inducing the map.  Decomposing this cycle via the K\"unneth
formula yields the K\"unneth projectors $\pi_i$, which is precisely
Conjecture~C.
\end{proof}

\begin{example}[Conjecture~A for curves]
For a smooth projective curve $X$ ($d=1$), we need $\mathbf{L}^0 = \id$ to
be algebraic, which is trivially true.  So Conjecture~A holds for curves.
For surfaces ($d=2$), Conjecture~A reduces to the statement that the
Lefschetz involution $\Lambda: H^2(X) \to H^2(X)$ is algebraic, which is
non-trivial.
\end{example}

% -------------------------------------------------------
\section{Conjecture~B: The Hodge Standard Conjecture}\label{sec:conjB}
% -------------------------------------------------------

Conjecture~B is in many ways the most fundamental of the four, as it implies
positivity properties of the intersection form.

\begin{definition}[Hodge-Riemann bilinear form]
Assuming hard Lefschetz (Conjecture~B), define for primitive classes
$\alpha, \beta \in H^i_{\mathrm{prim}}(X)$:
\[
  Q(\alpha, \beta) = (-1)^{i(i-1)/2} \int_X \alpha \cup \beta \cup L^{d-i}.
\]
This is the \emph{Hodge-Riemann bilinear form}.
\end{definition}

\begin{conjecture}[Standard Conjecture~B, Grothendieck 1969]\label{conj:B}
Let $X$, $H^*$, $d$, $\xi$ be as above.  Assume Conjecture~A holds for $X$.
Then:
\begin{enumerate}[label=(\roman*)]
  \item The Lefschetz involution $\Lambda: H^i(X) \to H^{i-2}(X)$ (for $i \geq 1$)
        is algebraic, i.e.\ induced by an algebraic correspondence.
  \item The Hodge-Riemann bilinear form $Q$ is \emph{positive definite} on
        each graded piece of the primitive cohomology.
\end{enumerate}
Equivalently: the $\mathfrak{sl}_2$-action generated by $\mathbf{L}$ and $\Lambda$
on $H^*(X)$ is algebraic.
\end{conjecture}

\begin{remark}
Over $\C$, Conjecture~B is a consequence of classical Hodge theory (the
Hodge-Riemann bilinear relations).  Conjecture~B over arbitrary fields is
entirely open for varieties of dimension $\geq 2$ (except in special cases,
cf.\ Section~\ref{sec:known}).
\end{remark}

\begin{proposition}[Conjecture~B $\Rightarrow$ Conjecture~A]
Standard Conjecture~B implies Standard Conjecture~A.
\end{proposition}

\begin{proof}
The algebraicity of $\Lambda$ and the formula $[\mathbf{L}, \Lambda] =
(d - i)\cdot\id$ on $H^i$ (the $\mathfrak{sl}_2$ relation) together force
$\mathbf{L}^{d-i}$ to be expressed in terms of $\Lambda$ and $\mathbf{L}$.
Since $\mathbf{L}$ is algebraic (cup product with a cycle class), and $\Lambda$
is algebraic by Conjecture~B, their composition is algebraic.
\end{proof}

\begin{theorem}[Grothendieck]\label{thm:B_implies_D}
Standard Conjecture~B (plus Conjecture~A) implies Standard Conjecture~D.
\end{theorem}

This is one of the deepest implications among the four conjectures.

% -------------------------------------------------------
\section{Conjecture~C: The K\"{u}nneth Conjecture}\label{sec:conjC}
% -------------------------------------------------------

Conjecture~C concerns the algebraicity of the K\"unneth components of the
diagonal.

\begin{definition}[K\"unneth decomposition of the diagonal]
By the K\"unneth formula for $H^*$, there is a decomposition
\[
  [\Delta_X] = \sum_{i=0}^{2d} \pi_i \in \bigoplus_{i=0}^{2d} H^{2d-i}(X)
  \otimes H^i(X) = H^{2d}(X \times X),
\]
where $[\Delta_X]$ is the class of the diagonal $\Delta_X \subset X \times X$
and $\pi_i$ is the component in $H^{2d-i}(X) \otimes H^i(X)$.
\end{definition}

\begin{conjecture}[Standard Conjecture~C, Grothendieck 1969]\label{conj:C}
For each $0 \leq i \leq 2d$, the K\"unneth component $\pi_i$ of the
diagonal class $[\Delta_X] \in H^{2d}(X \times X)(d)$ is algebraic, i.e.\
\[
  \pi_i \in \mathrm{Im}(\cl: \CH^d(X \times X) \otimes_\Z \Q \to H^{2d}(X \times X)).
\]
\end{conjecture}

\begin{remark}
Conjecture~C is equivalent to the existence of \emph{algebraic} projectors
$\pi_i \in \CH^d(X \times X)_\Q$ such that $H^*(X) = \bigoplus_i \pi_i H^*(X)$.
This is a necessary condition for $X$ to have a \emph{Chow-K\"unneth
decomposition} in the sense of Murre.
\end{remark}

\begin{proposition}[Properties of Conjecture~C]
\begin{enumerate}[label=(\roman*)]
  \item Conjecture~C is implied by Conjecture~A (as noted above).
  \item Conjecture~C implies that $\CH^*(X)_\Q$ decomposes as a direct sum
        of pieces $\CH^i(X)_\Q = \bigoplus_j \CH^i_j(X)_\Q$ (Bloch-Srinivas
        decomposition conjecture, in special cases).
  \item Over $\C$, Conjecture~C is known for all smooth projective varieties
        (by the existence of the Hodge decomposition, which gives explicit
        K\"unneth projectors as integrals of differential forms --- but these
        need not be algebraic in general!).
\end{enumerate}
\end{proposition}

\begin{remark}
The subtlety in point (iii) is that the Hodge decomposition gives K\"unneth
components in $H^{p,q}$, but these are defined using the complex structure and
transcendental methods.  Conjecture~C asks whether the \emph{same} projectors
can be represented by algebraic cycles.  Over $\C$ this is essentially the
Hodge conjecture applied to $X \times X$.
\end{remark}

% -------------------------------------------------------
\section{Conjecture~D: Numerical Equals Homological}\label{sec:conjD}
% -------------------------------------------------------

Conjecture~D is perhaps the most ``down to earth'' of the four, as it concerns
a purely cycle-theoretic statement that does not mention a specific cohomology
theory.

\begin{conjecture}[Standard Conjecture~D, Grothendieck 1969]\label{conj:D}
For any smooth projective variety $X$ over $k$ and any adequate equivalence
relation $\sim$ between homological and numerical equivalence (i.e.\ for any
prime $\ell \neq \ch(k)$ and the $\ell$-adic cohomology):
\[
  Z \sim_{\mathrm{hom}} 0 \iff Z \sim_{\mathrm{num}} 0,
  \quad \forall Z \in \mathcal{Z}^p(X).
\]
Equivalently, the cycle class map induces an injection
\[
  \cl: \CH^p(X)_\Q / {\sim_{\mathrm{num}}} \hookrightarrow H^{2p}(X)(p).
\]
\end{conjecture}

\begin{remark}
The inclusion ${\sim_{\mathrm{hom}}} \supseteq {\sim_{\mathrm{num}}}$ is
elementary (Poincar\'e duality).  The non-trivial assertion is the reverse:
if a cycle has zero intersection number with every cycle of complementary
dimension, then its cohomology class vanishes.
\end{remark}

\begin{theorem}[Kleiman 1968]\label{thm:D_implies_A}
Standard Conjecture~D implies Standard Conjecture~A.
\end{theorem}

\begin{proof}[Proof sketch]
Kleiman showed that if the cycle class map is injective on numerical
equivalence classes, then the K\"unneth decomposition of the diagonal can
be performed algebraically.  The key step is that the K\"unneth projector
$\pi_i$ is numerically trivial when combined with $\pi_j$ for $j \neq i$,
and Conjecture~D then forces it to be homologically trivial on the complement,
hence algebraic.  Full details appear in \cite{kleiman1968}.
\end{proof}

\begin{theorem}[Independence of Weil cohomology, conditional]\label{thm:independence}
Assuming Conjecture~D, the groups $\CH^p(X)_\Q/{\sim_{\mathrm{hom}}}$ are
independent of the choice of Weil cohomology theory $H^*$.
\end{theorem}

\begin{proof}
Conjecture~D equates homological and numerical equivalence, and numerical
equivalence is defined intrinsically (without reference to $H^*$).
\end{proof}

% -------------------------------------------------------
\section{Implications Between the Conjectures}\label{sec:implications}
% -------------------------------------------------------

We now give a complete picture of the logical dependencies.

\begin{theorem}[Implication diagram]\label{thm:implications}
The following implications hold unconditionally:
\[
  \mathbf{B} \Rightarrow \mathbf{A} \Rightarrow \mathbf{C},
  \qquad
  \mathbf{D} \Rightarrow \mathbf{A},
  \qquad
  \mathbf{A} + \mathbf{B} \Rightarrow \mathbf{D}.
\]
In particular, $\mathbf{B}$ implies all four conjectures.
\end{theorem}

\begin{proof}[Summary of proofs]
\begin{enumerate}[label=(\roman*)]
  \item $\mathbf{B} \Rightarrow \mathbf{A}$: Algebraicity of $\Lambda$ plus
        the $\mathfrak{sl}_2$ relations give algebraicity of $\mathbf{L}^{d-i}$.
  \item $\mathbf{A} \Rightarrow \mathbf{C}$: The K\"unneth components $\pi_i$
        are expressed via the algebraic correspondences for $\mathbf{L}^{d-i}$.
  \item $\mathbf{D} \Rightarrow \mathbf{A}$: Kleiman's theorem
        (Theorem~\ref{thm:D_implies_A}).
  \item $\mathbf{A} + \mathbf{B} \Rightarrow \mathbf{D}$: Positive
        definiteness (from $\mathbf{B}$) plus algebraic projectors (from $\mathbf{A}$)
        force every homologically trivial cycle to be numerically trivial.
\end{enumerate}
\end{proof}

\begin{corollary}
If any one of the four standard conjectures is proved for \emph{all} smooth
projective varieties, then we get substantial progress toward the others.
Specifically: $\mathbf{B} \Rightarrow \mathbf{A}, \mathbf{C}, \mathbf{D}$
all simultaneously.
\end{corollary}

\begin{remark}[What is \emph{not} known]
It is not known whether $\mathbf{C}$ or $\mathbf{D}$ individually implies
$\mathbf{B}$.  The implication $\mathbf{D} \Rightarrow \mathbf{B}$ is open.
\end{remark}

% -------------------------------------------------------
\section{Relation to the Weil Conjectures}\label{sec:weilconj}
% -------------------------------------------------------

The Weil conjectures (proved by Deligne in 1974 \cite{deligne1974}) assert that
for a smooth projective variety $X$ over $\Fp$, the zeta function
\[
  Z(X, T) = \exp\!\left(\sum_{n=1}^\infty |X(\mathbb{F}_{p^n})| \frac{T^n}{n}\right)
\]
satisfies:
\begin{enumerate}[label=(\roman*)]
  \item \textbf{Rationality}: $Z(X,T) \in \Q(T)$.
  \item \textbf{Functional equation}: $Z(X, q^{-d}T^{-1}) = \pm q^{d\chi/2} T^\chi Z(X,T)$.
  \item \textbf{Riemann hypothesis}: The eigenvalues of Frobenius on $H^i$ have
        absolute value $p^{i/2}$.
  \item \textbf{Betti numbers}: $Z(X,T) = \prod_{i=0}^{2d} P_i(T)^{(-1)^{i+1}}$
        where $\deg P_i = b_i(X)$.
\end{enumerate}

\begin{theorem}[Grothendieck's conditional proof]
Assuming the standard conjectures (specifically Conjecture~B), one can prove
the Riemann hypothesis part of the Weil conjectures.
\end{theorem}

\begin{remark}
This was Grothendieck's original motivation for formulating the standard
conjectures.  However, Deligne found a completely different proof of the
Weil conjectures that \emph{does not} use the standard conjectures.  This
makes the standard conjectures even more mysterious: they imply more about
the structure of cycles and motives than is needed merely for counting points.
\end{remark}

\begin{proposition}
The Weil conjectures (proved) do \emph{not} imply any of the standard
conjectures.  The standard conjectures are strictly stronger.
\end{proposition}

% -------------------------------------------------------
\section{Relation to the Hodge and Tate Conjectures}\label{sec:hodge}
% -------------------------------------------------------

\begin{conjecture}[Hodge conjecture, Hodge 1950]
Let $X$ be a smooth projective complex variety.  Then every rational Hodge
class in $H^{2p}(X,\Q) \cap H^{p,p}(X)$ is a rational linear combination of
cycle classes $\cl(Z)$ for algebraic cycles $Z \in \mathcal{Z}^p(X)$.
\end{conjecture}

\begin{proposition}[Hodge implies A over $\mathbb{C}$]
The Hodge conjecture (for $X$ and $X \times X$) implies Standard Conjecture~A
over $\C$.
\end{proposition}

\begin{proof}
The cycle class $[\Delta_X] \in H^{2d}(X \times X, \Q)$ has K\"unneth
components $\pi_i \in H^{2d-i}(X) \otimes H^i(X)$.  By classical Hodge theory,
each $\pi_i$ is a Hodge class (it lies in $H^{d,d}(X \times X)$ when
$i$ ranges over even degrees in the appropriate sense).  The Hodge conjecture
applied to $X \times X$ then gives algebraicity of $\pi_i$.
\end{proof}

\begin{conjecture}[Tate conjecture, Tate 1963]
Let $X$ be a smooth projective variety over a finitely generated field $k$,
and $\ell$ a prime different from $\ch(k)$.  Then the cycle class map
\[
  \cl: \CH^p(X)_{\Q_\ell} \to H^{2p}(X_{\bar k}, \Q_\ell(p))^{G_k}
\]
is surjective (where $G_k = \Gal(\bar k/k)$).
\end{conjecture}

\begin{theorem}[Tate implies D over finite fields]
The Tate conjecture (surjectivity) plus its injectivity variant implies
Standard Conjecture~D over finite fields.
\end{theorem}

The precise relation between the Tate conjecture and the standard conjectures
over finite fields was worked out by Tate~\cite{tate1994} and Milne~\cite{milne2002}.

% -------------------------------------------------------
\section{Motivic Consequences}\label{sec:motives}
% -------------------------------------------------------

The deepest justification for the standard conjectures comes from their role
in the theory of motives.

\begin{definition}[Category of pure motives]
Fix an adequate equivalence relation $\sim$ on algebraic cycles.  The
\emph{category of pure motives} $\mathcal{M}_\sim(k)$ has:
\begin{itemize}
  \item \textbf{Objects}: pairs $(X, p)$ where $X$ is smooth projective and
        $p \in \CH^d(X \times X)_\Q$ is an idempotent correspondence (projector).
  \item \textbf{Morphisms}: $\Hom((X,p),(Y,q)) = q \circ \CH^d(X \times Y)_\Q \circ p$.
  \item \textbf{Tensor product}: $(X,p) \otimes (Y,q) = (X \times Y, p \times q)$.
\end{itemize}
\end{definition}

\begin{theorem}[Jannsen 1992]\label{thm:jannsen}
The category $\mathcal{M}_{\mathrm{num}}(k)$ of pure motives modulo numerical
equivalence is a \emph{semi-simple abelian category}.
\end{theorem}

This theorem, proved by Jannsen without assuming any standard conjectures, is
one of the major unconditional results in motivic theory.

\begin{theorem}[Conditional structure theorem for motives]
Assuming the standard conjectures, the functor
\[
  h: \{\text{smooth projective }k\text{-varieties}\} \to \mathcal{M}_{\mathrm{num}}(k)
\]
$X \mapsto (X, [\Delta_X])$ decomposes as
$h(X) = \bigoplus_{i=0}^{2d} h^i(X)$,
where each $h^i(X)$ is a pure motive of weight $i$.  Moreover,
$\mathcal{M}_{\mathrm{num}}(k)$ is a Tannakian category and satisfies
axioms analogous to those of a pure Hodge structure category.
\end{theorem}

\begin{remark}
The Tannakian structure would allow one to define a ``motivic Galois group''
$G_{\mathrm{mot}}(k)$ whose representations classify the motives.  The standard
conjectures are necessary to make this group well-defined and algebraic.
\end{remark}

% -------------------------------------------------------
\section{Andr\'{e}'s Motivated Cycles}\label{sec:andre}
% -------------------------------------------------------

In the absence of a proof of the standard conjectures, Andr\'e~\cite{andre1996}
introduced a larger class of ``motivated cycles'' that satisfies all the
desired algebraic properties.

\begin{definition}[Motivated cycles, Andr\'{e} 1996]
A cohomology class $\alpha \in H^{2p}(X)(p)$ is called a \emph{motivated
cycle} if it belongs to the smallest $\Q$-subalgebra of $\bigoplus_X H^*(X)$
generated by:
\begin{enumerate}[label=(\roman*)]
  \item Algebraic cycle classes $\cl(\mathcal{Z}^p(X))$.
  \item Classes of the form $*_L(\beta)$ for motivated $\beta$, where $*_L$
        is the Hodge-$*$ operator associated to some polarisation $L$.
\end{enumerate}
\end{definition}

\begin{theorem}[Andr\'{e} 1996]\label{thm:andre}
The category $\mathcal{M}_{\mathrm{mot}}(k)$ of motives built from motivated
cycles is:
\begin{enumerate}[label=(\roman*)]
  \item A semi-simple $\Q$-linear abelian tensor category.
  \item Satisfies all four standard conjectures (by construction).
  \item Contains the abelian motives as a subcategory.
\end{enumerate}
\end{theorem}

\begin{remark}
Andr\'e's theorem does not prove the standard conjectures: it only shows that
\emph{if} one enlarges the class of algebraic cycles to include the $*_L$
operator, then the resulting category behaves as expected.  The standard
conjectures ask whether $*_L$ is already algebraic.
\end{remark}

\begin{theorem}[Motivated cycles and the Lefschetz group, Andr\'{e}]
The Lefschetz group $G_L(X)$ (the Zariski closure of the monodromy) acts on
motivated cycles, and the motivated cohomology is stable under this action.
This gives a precise sense in which motivated cycles behave ``like''
algebraic cycles in classical Hodge theory.
\end{theorem}

% -------------------------------------------------------
\section{Jannsen's Semisimplicity Theorem}\label{sec:jannsen}
% -------------------------------------------------------

\begin{theorem}[Jannsen 1992, full statement]\label{thm:jannsen_full}
Let $k$ be any field.  The category $\mathcal{M}_{\mathrm{num}}(k)$ of pure
motives modulo numerical equivalence (over $k$, with $\Q$-coefficients) is a
semi-simple abelian rigid tensor category.
\end{theorem}

\begin{proof}[Proof sketch]
The key steps are:
\begin{enumerate}[label=(\arabic*)]
  \item Show that $\End_{\mathcal{M}_{\mathrm{num}}}(h(X))$ is a
        finite-dimensional $\Q$-algebra for all $X$.
  \item Show this algebra has no nilpotents (using the fact that numerical
        equivalence is the ``coarsest'' adequate equivalence relation, and
        that $\CH^*(X)_\Q/{\sim_{\mathrm{num}}}$ is a $\Q$-algebra with positive
        definite intersection form on each degree).
  \item Conclude that $\mathcal{M}_{\mathrm{num}}(k)$ is semi-simple by a
        general categorical argument.
\end{enumerate}
\end{proof}

\begin{corollary}
Every motive in $\mathcal{M}_{\mathrm{num}}(k)$ is a direct sum of simple
motives.  If the standard conjectures hold, each simple motive corresponds
to an irreducible Galois representation.
\end{corollary}

% -------------------------------------------------------
\section{Known Results and Special Cases}\label{sec:known}
% -------------------------------------------------------

\begin{theorem}[Known cases of Conjecture~A]
Standard Conjecture~A is known for:
\begin{enumerate}[label=(\roman*)]
  \item Smooth projective curves ($d=1$): trivial.
  \item Smooth projective surfaces ($d=2$): follows from the algebraicity of
        the Lefschetz involution $\Lambda$ for divisors, which is Lefschetz's
        theorem on $(1,1)$-classes (Lefschetz~\cite{lefschetz1924}).
  \item Abelian varieties: Lieberman~\cite{lieberman1968} proved Conjecture~A
        for abelian varieties using the theory of the Rosati involution.
  \item Varieties with a cellular decomposition (e.g.\ flag varieties,
        Grassmannians, toric varieties): K\"unneth components are algebraic
        by explicit construction.
\end{enumerate}
\end{theorem}

\begin{theorem}[Known cases of Conjecture~B]
Standard Conjecture~B (Hodge standard conjecture) is known for:
\begin{enumerate}[label=(\roman*)]
  \item Curves and surfaces over $\C$: classical Hodge theory.
  \item Abelian varieties over any field: follows from the Rosati involution
        (see Mumford~\cite{mumford1970}).
  \item $K3$ surfaces: verified by Kleiman and others.
\end{enumerate}
Conjecture~B is \emph{open} for all varieties of dimension $\geq 3$ over fields
of positive characteristic.
\end{theorem}

\begin{theorem}[Known cases of Conjecture~D]
Conjecture~D is known for:
\begin{enumerate}[label=(\roman*)]
  \item Divisors on any smooth projective variety: this is a theorem
        (follows from the fact that numerical and homological equivalence agree
        for $\CH^1$; the latter is the N\'eron-Severi theorem).
  \item Zero-cycles on surfaces: by a theorem of Bloch and Srinivas~\cite{bloch1983}.
  \item Abelian varieties: known by the work of Lieberman and others.
\end{enumerate}
\end{theorem}

% -------------------------------------------------------
\section{Obstacles to a General Proof}\label{sec:obstacles}
% -------------------------------------------------------

\begin{remark}[Why the standard conjectures are hard]
The fundamental difficulty is the \emph{transcendence} of the Hodge star
operator $*_L$.  In classical Hodge theory over $\C$, $*_L$ is defined using
the metric on the variety, which has no algebraic interpretation.  The
standard conjectures ask for a purely algebraic substitute.

More precisely:
\begin{enumerate}[label=(\roman*)]
  \item There is no known algebraic construction of $\Lambda$ that works in
        positive characteristic.
  \item The Witt cohomology and crystalline cohomology approaches have not
        yet yielded a proof of Conjecture~B.
  \item The Langlands programme gives no direct handle on cycles.
  \item Motives over finite fields are better understood (via Frobenius), but
        even there Conjecture~B remains open for $d \geq 3$.
\end{enumerate}
\end{remark}

\begin{remark}[The p-adic approach]
Berthelot's crystalline cohomology and Fontaine's $p$-adic Hodge theory have
been used to prove special cases.  In particular, for varieties with good
reduction at $p$, the Katz-Messing theorem gives the Weil conjectures from
Hodge theory, but it does not prove the standard conjectures.
\end{remark}

% -------------------------------------------------------
\section{Outlook and Open Problems}\label{sec:outlook}
% -------------------------------------------------------

\begin{conjecture}[Murre's conjectures, 1993]
For every smooth projective variety $X$ of dimension $d$, there exist
idempotents $\Pi_0, \Pi_1, \ldots, \Pi_{2d} \in \CH^d(X \times X)_\Q$ such
that $\sum_i \Pi_i = [\Delta_X]$, $\Pi_i \circ \Pi_j = \delta_{ij} \Pi_i$,
and $(\Pi_i)_*$ acts on $H^j(X)$ as $\delta_{ij}\cdot\id$.
\end{conjecture}

This is a strong form of Conjecture~C (with the explicit idempotents satisfying
additional orthogonality conditions).

\begin{remark}[The motivic cohomology programme]
Voevodsky's theory of motivic cohomology~\cite{voevodsky2000} provides a
framework for understanding algebraic cycles that goes beyond the standard
conjectures.  In particular, the Milnor conjecture (proved by Voevodsky) and
the Bloch-Kato conjecture (proved by Voevodsky-Rost) give deep results about
the structure of Milnor $K$-theory, but do not directly imply the standard
conjectures.
\end{remark}

\begin{remark}[Potential approaches via derived categories]
Bondal and Orlov's reconstruction theorem~\cite{bondal2001} suggests that
the derived category $D^b(\mathrm{Coh}(X))$ contains all motivic information.
If one could relate derived equivalences to motivated cycles, this might give
a new approach to the standard conjectures.
\end{remark}

\begin{remark}[The ``mixed'' case]
The standard conjectures are formulated for \emph{pure} motives (smooth
projective varieties).  For \emph{mixed} motives (general varieties), the
corresponding conjectures become even harder.  Beilinson's conjectures
\cite{beilinson1985} generalise the Bloch-Kato framework to this setting.
\end{remark}

\begin{thebibliography}{99}

\bibitem{grothendieck1968}
A.\ Grothendieck,
\textit{Standard conjectures on algebraic cycles},
in: Algebraic Geometry (Internat.\ Colloq., Tata Inst.\ Fund.\ Res., Bombay, 1968),
Oxford Univ.\ Press, London, 1969, pp.\ 193--199.

\bibitem{kleiman1968}
S.\ L.\ Kleiman,
\textit{Algebraic cycles and the Weil conjectures},
in: Dix expos\'es sur la cohomologie des sch\'emas, Masson \& Cie, Paris; North-Holland, Amsterdam, 1968, pp.\ 359--386.

\bibitem{deligne1974}
P.\ Deligne,
\textit{La conjecture de Weil, I},
Publ.\ Math.\ IHES \textbf{43} (1974), 273--307.

\bibitem{deligne1980}
P.\ Deligne,
\textit{La conjecture de Weil, II},
Publ.\ Math.\ IHES \textbf{52} (1980), 137--252.

\bibitem{jannsen1992}
U.\ Jannsen,
\textit{Motives, numerical equivalence, and semi-simplicity},
Invent.\ Math.\ \textbf{107} (1992), 447--452.

\bibitem{andre1996}
Y.\ Andr\'e,
\textit{Pour une th\'eorie inconditionnelle des motifs},
Publ.\ Math.\ IHES \textbf{83} (1996), 5--49.

\bibitem{tate1994}
J.\ Tate,
\textit{Conjectures on algebraic cycles in $\ell$-adic cohomology},
in: Motives (Seattle, 1991), Proc.\ Sympos.\ Pure Math.\ \textbf{55}, Part 1,
Amer.\ Math.\ Soc., Providence, RI, 1994, pp.\ 71--83.

\bibitem{milne2002}
J.\ S.\ Milne,
\textit{Lefschetz motives and the Tate conjecture},
Compositio Math.\ \textbf{117} (1999), 45--76.

\bibitem{lieberman1968}
D.\ Lieberman,
\textit{Numerical and homological equivalence of algebraic cycles on Hodge manifolds},
Amer.\ J.\ Math.\ \textbf{90} (1968), 366--374.

\bibitem{mumford1970}
D.\ Mumford,
\textit{Abelian Varieties},
Tata Institute of Fundamental Research Studies in Mathematics, Oxford Univ.\ Press, 1970.

\bibitem{bloch1983}
S.\ Bloch, V.\ Srinivas,
\textit{Remarks on correspondences and algebraic cycles},
Amer.\ J.\ Math.\ \textbf{105} (1983), 1235--1253.

\bibitem{lefschetz1924}
S.\ Lefschetz,
\textit{L'Analysis Situs et la G\'eom\'etrie Alg\'ebrique},
Gauthier-Villars, Paris, 1924.

\bibitem{voevodsky2000}
V.\ Voevodsky,
\textit{Triangulated categories of motives over a field},
in: Cycles, Transfers, and Motivic Homology Theories,
Ann.\ of Math.\ Stud.\ \textbf{143}, Princeton Univ.\ Press, 2000, pp.\ 188--238.

\bibitem{bondal2001}
A.\ Bondal, D.\ Orlov,
\textit{Reconstruction of a variety from the derived category and groups of autoequivalences},
Compositio Math.\ \textbf{125} (2001), 327--344.

\bibitem{beilinson1985}
A.\ A.\ Beilinson,
\textit{Higher regulators and values of $L$-functions},
J.\ Soviet Math.\ \textbf{30} (1985), 2036--2070.

\bibitem{kleiman1994}
S.\ L.\ Kleiman,
\textit{The standard conjectures},
in: Motives (Seattle, 1991), Proc.\ Sympos.\ Pure Math.\ \textbf{55}, Part 1,
Amer.\ Math.\ Soc., Providence, RI, 1994, pp.\ 3--20.

\bibitem{murre1993}
J.\ P.\ Murre,
\textit{On a conjectural filtration on the Chow groups of an algebraic variety, I, II},
Indag.\ Math.\ (N.S.) \textbf{4} (1993), 177--201.

\bibitem{scholl1994}
A.\ J.\ Scholl,
\textit{Classical motives},
in: Motives (Seattle, 1991), Proc.\ Sympos.\ Pure Math.\ \textbf{55}, Part 1,
Amer.\ Math.\ Soc., Providence, RI, 1994, pp.\ 163--187.

\end{thebibliography}

\end{document}
